\chapter{Conceptos básicos de estabilidad y estabilidad para sistemas lineales}
Estabilidad : si cambiando ligeramente las condiciones iniciales la solución cambia poco para todo tiempo $\displaystyle t \geq t_{0} $. 
\section{Conceptos de estabilidad}
Supongamos que $\displaystyle f : D \subset \R \times \R^{n} \to \R^{n} $ es una función continua y localmente Lipschitz respecto a la segunda variable. Sea $\displaystyle \overline{x} : [\tau, \infty) \to \R $ que satisface $\displaystyle \overline{x}' = f\left(t, \overline{x}\right) $, $\displaystyle \forall t \geq \tau $. 
\begin{definition}[Estabilidad]
Diremos que $\displaystyle \overline{x} $ es una solución \textbf{estable} si para todo $\displaystyle \epsilon > 0 $ y $\displaystyle t_{0}\geq \tau $ existe $\displaystyle \delta > 0 $ tal que si $\displaystyle \|x_{0}-\overline{x}\left(t_{0}\right)\| < \delta  $ entonces $\displaystyle \|\varphi\left(t; t_{0}, x_{0}\right)-\overline{x}\left(t\right)\| < \epsilon  $, $\displaystyle \forall t \geq t_{0} $. Diremos que una solución es \textbf{inestable} si no es estable.  
\end{definition}
\begin{observation}
Para que una solución sea estable tiene que estar definida en un intervalo $\displaystyle [\tau, \infty )$ para cierto $\displaystyle \tau \in [-\infty, \infty) $. 
\end{observation}
\begin{definition}[Estabilidad asintótica]
Diremos que $\displaystyle \overline{x} $ es una solución \textbf{asintóticamente estable} si es estable y para todo $\displaystyle t_{0}\geq \tau $ eixste $\displaystyle \eta > 0 $ tal que si $\displaystyle \left|x_{0}-\overline{x}\left(t_{0}\right)\right|<\eta  $ entonces
\[ \left|\varphi\left(t; t_{0}, x_{0}\right)-\overline{x}\left(t_{0}\right)\right|\xrightarrow{t \to \infty} 0 .\]
\end{definition}
\begin{eg}
Consideremos la EDO $\displaystyle x' = ax $. Tendremos que 
\[\varphi\left(t; t_{0}, x_{0}\right)=x_{0}e^{a\left(t-t_{0}\right)} .\]
Distinguimos varios casos:
\begin{itemize}
\item Si $\displaystyle a \leq 0 $, cogemos $\displaystyle \overline{x}\left(t\right)=\overline{x}_{0}e^{a\left(t-\tau\right)} $. Entonces, 
	\[
	\begin{split}
		\left|\varphi\left(t; t_{0}, x_{0}\right)-\overline{x}\left(t\right)\right|= & \left|x_{0}e^{a\left(t-t_{0}\right)}-\overline{x}_{0}e^{a\left(t-\tau\right)}\right| \\
		= & e^{a\left(t-t_{0}\right)} \left|x_{0}-\overline{x}_{0}e^{t_{0}-\tau}\right| \leq \left|x_{0}-\overline{x}\left(t\right)\right|<\epsilon  .
	\end{split}
	\]
Basta con coger $\displaystyle \delta = \epsilon  $ para comprobar la estabilidad. 
Estudiemos ahora la estabilidad asimptótica. Si $\displaystyle a > 0 $, tenemos que 
	\[ \left|\varphi\left(t; t_{0}, x_{0}\right)-\overline{x}\left(t\right)\right|= \left|x_{0}-\overline{x}\left(t_{0}\right)\right|e^{a\left(t-t_{0}\right)} .\]
Podemos elegir $\displaystyle \eta > 0 $ y obtenemos la estabilidad asintótica. Si $\displaystyle a = 0 $ se puede ver que que es estable pero no asintóticamente estable. 
\item Si $\displaystyle a > 0 $, sea cual sea $\displaystyle \delta  $, si $\displaystyle \left|x_{0}-\overline{x}\left(t_{0}\right)\right| < \delta  $ tenemos que 
	\[ \left|\varphi\left(t; t_{0}, x_{0}\right)-\overline{x}\left(t_{0}\right)\right|= \left|x_{0}-\overline{x}\left(t_{0}\right)\right|e^{a\left(t-t_{0}\right)} \to \infty ,\]
por tanto es una solución inestable. 	
\end{itemize}
\end{eg}
\begin{eg}
Consideremos la EDO $\displaystyle x ' = x^{2} $, podemos coger
\[\varphi\left(t; t_{0}, x_{0}\right)=\frac{x_{0}}{1-x_{0}\left(t-t_{0}\right)} .\]
Otro ejemplo es $\displaystyle x' = x\left(1-x\right) $.
\end{eg}
\begin{lema}[Reducción de estabilidad a la trivial]
Dado el sistema $\displaystyle x' = f\left(t,x\right) $ y una solución $\displaystyle \overline{x} : [\tau, \infty) \to \R^{n} $. Si consideramos la función
\[g\left(t,y\right)=f\left(t,y + \overline{x}\left(t\right)\right)-f\left(t, \overline{x}\left(t\right)\right) ,\]
entonces
\begin{itemize}
\item $\displaystyle \overline{y}\left(t\right)\equiv 0 $ es solución de $\displaystyle y' = g\left(t,y\right) $. 
\item $\displaystyle \overline{x} $ es solución estable (resp. asintóticamente estable) de $\displaystyle x' = f\left(t,x\right) $ si y sólo si $\displaystyle \overline{y} $ es solución estable (resp. asintóticamente estable) de $\displaystyle y' = g\left(t,y\right) $. 
\end{itemize}
\end{lema}
\begin{proof}
En primer lugar, tenemos que 
\[g\left(t, 0\right)=f\left(t, \overline{x}\left(t\right)\right)-f\left(t, \overline{x}\left(t\right)\right)=0 = \overline{y}'\left(t\right) ,\]
habiendo tomado $\displaystyle \overline{y}\left(t\right)=0 $. Ahora, sea $\displaystyle \varphi\left(t; t_{0}, x_{0}\right) $ el flujo de $\displaystyle x' = f\left(t,x\right) $ y $\displaystyle \phi\left(t; t_{0}, x_{0}\right) $ el flujo de $\displaystyle y' = g\left(t,y\right) $. Entonces tenemos que, 
\[\phi\left(t; t_{0}, y_{0}\right)=\varphi\left(t, t_{0}, y_{0} + \overline{x}\left(t_{0}\right)\right)-\overline{x}\left(t_{0}\right) .\]
Si cogemos $\displaystyle z\left(t\right):=\varphi\left(t, t_{0}, y_{0} + \overline{x}\left(t_{0}\right)\right)-\overline{x}\left(t\right)  $, tenemos que 
\[
\begin{split}
	z'\left(t\right)= &\frac{d}{dt}\left(\varphi\left(t, t_{0}, y_{0} + \overline{x}\left(t_{0}\right)\right)-\overline{x}\left(t\right) \right)=\frac{d\varphi}{dt}\left(t; t_{0}, y_{0}+\overline{x}\left(t_{0}\right)\right)-\overline{x}'\left(t\right)\\
	= & f\left(t; \varphi\left(t; t_{0}, y_{0} + \overline{x}\left(t_{0}\right)\right)\right)-f\left(t, \overline{x}\left(t\right)\right) 
	=  f\left(t, z\left(t\right)+\overline{x}\left(t\right)\right)-f\left(t, \overline{x}\left(t\right)\right) \\
	= &  g\left(t, z\left(t\right)\right).
\end{split}
\]
Además, 
\[z\left(t_{0}\right)=\varphi\left(t_{0}; t_{0}, y_{0} + \overline{x}\left(t_{0}\right)\right)-\overline{x}\left(t_{0}\right)=y_{0} + \overline{x}\left(t_{0}\right)-\overline{x}\left(t_{0}\right)=y_{0} .\]
\end{proof}
\section{Estabilidad para sistemas lineales}
Consideremos el sistema 
\[x' = A\left(t\right)x + b\left(t\right) ,\]
con $\displaystyle A\in \mathcal{C}\left(I, \R^{n \times m}\right) $ y $\displaystyle b \in \mathcal{C}\left(I, \R^{n}\right) $, con $\displaystyle I $ un intervalo de $\displaystyle \R $. 
\begin{observation}
Con estas hipótesis de $\displaystyle A $ y $\displaystyle B $ podemos tomar
\[f\left(t,x\right)=A\left(t\right)x + b ,\]
y $\displaystyle f $ cumple las hipótesis del teorema Picard Lindelöf local. En efecto, 
\[ \|f\left(t,x\right)-f\left(t,y\right)\| = \|A\left(t\right)\left(x-y\right)\| \leq \|A\left(t\right)\| \|x-y\| .\]
\end{observation}
\begin{prop}
Todas las soluciones de un sistema lineal $\displaystyle x' = A\left(t\right)x + b\left(t\right) $ tienen la misma estabilidad y además coincide con la de la solución $\displaystyle y \equiv 0 $ del sistema homogéneo $\displaystyle y' = A\left(t\right)y $. 
\end{prop}
\begin{proof}
Usamos el lema anterior dada una solución $\displaystyle \overline{x}\left(t\right) $. Así, construimos 
\[
\begin{split}
	g\left(t,y\right)= & f\left(t,y+\overline{x}\left(t\right)\right)-f\left(t, \overline{x}\left(t\right)\right) \\
	= &  \left(A\left(t\right)\left(y + \overline{x}\left(t\right)+b\left(t\right)\right)\right)-\left(A\left(t\right)\overline{x}\left(t\right)+b\left(t\right)\right)=A\left(t\right)y.
\end{split}
\]
\end{proof}
\section*{Recordatorio}
El flujo del sistema $\displaystyle x' = A\left(t\right)x $ tiene la forma 
\[\varphi\left(t; t_{0}, x_{0}\right)=\Phi\left(t; t_{0}\right)x_{0} ,\]
donde $\displaystyle \Phi\left(t, t_{0}\right) $ es la matriz fundamental. Además, 
\[
\begin{cases}
\frac{d\Phi}{dt}\left(t; t_{0}\right)=A\left(t\right)\Phi\left(t; t_{0}\right) \\
\Phi\left(t_{0}; t_{0}\right)=I
\end{cases}
.\]
Por la propiedad del semigrupo tenemos que $\displaystyle \varphi\left(t; t_{0}, x_{0}\right)=\varphi\left(t; s, \varphi\left(s; t_{0}, x_{0}\right)\right)  $ por lo que
\[ \Phi\left(t; t_{0}\right)x_{0}=\Phi\left(t,s\right)\varphi\left(s; t_{0}, x_{0}\right)=\Phi\left(t,s\right)\Phi\left(s;t_{0}\right)x_{0}.\]
Como esto es cierto para todo $\displaystyle x_{0}\in \R^{n} $, debe ser que $\displaystyle \Phi\left(t; t_{0}\right)=\Phi\left(t; s\right)\Phi\left(s; t_{0}\right) $. En particular, si $\displaystyle t = t_{0} $,
\[I = \Phi\left(t_{0}; t_{0}\right)=\Phi\left(t_{0}, s\right)\Phi\left(s; t_{0}\right) .\]
\begin{theorem}[Criterio de estabilidad para sistemas lineales homogéneos]
Un sistema lineal homogéneo $\displaystyle y' = A\left(t\right)y $, es estable si y sólo si todas sus soluciones son acotadas. 
\end{theorem}
\begin{proof}
Consideramos ambas implicaciones.
\begin{description}
\item[($\displaystyle \Rightarrow $)] Como el sistema es estable, tenemos que para cada $\displaystyle \epsilon > 0 $ y para cada $\displaystyle t_{0} \in \R $, existe $\displaystyle \delta > 0 $ tal que si $\displaystyle \|x_{0}\| < \delta  $ entonces $\displaystyle \|\varphi\left(t; t_{0}, x_{0}\right)\|< \epsilon  $. Estamos cogiendo $\displaystyle \overline{x}\equiv 0 $. Elegimos $\displaystyle \epsilon = 1 $, por tanto existe $\displaystyle \delta > 0 $ tal que si $\displaystyle \|x_{0}\| < \delta  $ entonces $\displaystyle \|\varphi\left(t; t_{0}, x_{0}\right)\| < 1 $. 
	Si $\displaystyle x_{0}\in \R^{n} $,
	\[
	\begin{split}
		\|\varphi\left(t; t_{0}, x_{0}\right)\| = & \|\Phi\left(t; t_{0}, x_{0}\right)\| = \left\|\Phi\left(t; t_{0}\right)\frac{x_{0}}{\|x_{0}\|}\frac{\delta }{2}\right\|\frac{2\|x_{0}\|}{\delta } \\
		= &  \left\|\varphi\left(t; t_{0}, \frac{x_{0}\delta }{2\|x_{0}\|}\right)\right\|\frac{2\|x_{0}\|}{\delta } < \frac{2\|x_{0}\|}{\delta }.
	\end{split}
	\]
	Así, la solución está acotada. 
\item[($\displaystyle \Leftarrow $)] Sabemos que $\displaystyle \|\varphi\left(t; t_{0}, x_{0}\right)\| $ está acotada. En particular, la matriz fundamental está acotada. Así, $\displaystyle \|\Phi\left(t, t_{0}\right)\| \leq C $, $\displaystyle \forall t \geq t_{0} $. Entonces, $\displaystyle \forall \epsilon > 0 $ cogemos $\displaystyle \delta = \frac{\epsilon }{C} $ y hacemos un proceso como el anterior. 
\end{description}
\end{proof}

